\documentclass[sigconf]{acmart}

% CIKM Demo papers are short; keep the layout tight but compliant.
\settopmatter{authorsperrow=3}
\settopmatter{printacmref=false}
\settopmatter{printccs=false}
\settopmatter{printfolios=true}

\AtBeginDocument{%
  \providecommand\BibTeX{{Bib\TeX}}}

% Rights / metadata (update for CIKM'26) , -
\setcopyright{acmlicensed}
\copyrightyear{2026}
\acmYear{2026}
\acmDOI{XXXXXXX.XXXXXXX}

\acmConference[CIKM '26]{Proceedings of the 35th ACM International Conference on Information and Knowledge Management}{October 19--23, 2026}{Mexico City, Mexico}
\acmISBN{978-1-4503-XXXX-X/26/10}

% Packages (keep minimal but demo-ready) , -
\usepackage{booktabs}
\usepackage{url}
\usepackage{enumitem}
\usepackage{graphicx}
\usepackage{subcaption}
\usepackage{microtype}

\begin{document}

\title{GraphCraft: Compiling a Curator's Edit Journal into Reusable Skills with End-to-End Provenance}

\settopmatter{authorsperrow=3}

% Authors
\author{Ilia Karpov}
\orcid{0000-0002-8106-9426}
\email{karpovilia@gmail.com}
\affiliation{%
  \institution{HSE University, APSolutions LLC}
  \city{Moscow}
  \country{Russia}}

\author{Ivan Chepikov}
\orcid{0009-0004-4256-5437}
\email{ivanchepikov444@gmail.com}
\affiliation{%
  \institution{HSE University, APSolutions LLC}
  \city{Moscow}
  \country{Russia}}

\author{Denis Losev}
\email{denislosev48@gmail.com}
\affiliation{%
  \institution{APSolutions LLC}
  \city{Moscow}
  \country{Russia}}

\author{Alina Avanesyan}
\email{aavanesyan@hse.ru}
\affiliation{%
  \institution{HSE University}
  \city{Moscow}
  \country{Russia}}

\author{Mary Godunova}
\email{mgodunova@hse.ru}
\affiliation{%
  \institution{HSE University}
  \city{Moscow}
  \country{Russia}}

\author{Dmitry Ilvovsky}
\orcid{0000-0002-5484-372X}
\email{dilvovsky@hse.ru}
\affiliation{%
  \institution{HSE University}
  \city{Moscow}
  \country{Russia}}

\renewcommand{\shortauthors}{Karpov et al.}


\begin{abstract}
GraphRAG pipelines build entity-relation graphs from text and use them for retrieval and question answering, but in practice they operate as a black box: extraction noise, fragmented entities, and unstable clustering silently distort the evidence behind an answer, and fixing them is manual, repetitive work that no two corpora share, because domain, purpose, and language dictate which errors dominate.

We present GraphCraft, an interactive curation system built around two ideas. First, the journal of the curator's edits is training data: any slice of it can be compiled into a reusable skill. An agent reads the last $n$ operations as structured witnesses, induces the rule the curator follows, and a wizard materializes that rule at three price points, from a zero-LLM heuristic to a fully agentic build; each build is dry-run on the live graph and adopted only after the user has compared their previewed effects and costs. Skills may contradict one another, so proposals never auto-apply, and a node that a human has verified can be pinned, excluding it from every agent. Second, nothing is asserted without provenance: every answer traces to source spans through its evidence subgraph, every suggestion carries the rule and journal slice that produced it, and every automatic invalidation records its cause and can be reverted. The surrounding machinery, a bi-temporal store, a visual delta grammar, an activity timeline, incremental recomputation, and multi-user editing, is deliberately treated as supporting tooling. We demonstrate the workflow on a podcast corpus and an evolving news corpus and release the system as an open-source, offline-deployable platform.
\footnote{Code and offline deployment: \url{https://github.com/karpovilia/graphrag}; demo video: \url{https://vimeo.com/1141894235}.}
\end{abstract}

\begin{CCSXML}
<ccs2012>
  <concept>
    <concept_id>10002951.10003317</concept_id>
    <concept_desc>Information systems~Information retrieval</concept_desc>
    <concept_significance>500</concept_significance>
  </concept>
  <concept>
    <concept_id>10010147.10010178.10010179</concept_id>
    <concept_desc>Computing methodologies~Natural language processing</concept_desc>
    <concept_significance>500</concept_significance>
  </concept>
  <concept>
    <concept_id>10003120.10003145.10003147.10010365</concept_id>
    <concept_desc>Human-centered computing~Visual analytics</concept_desc>
    <concept_significance>300</concept_significance>
  </concept>
</ccs2012>
\end{CCSXML}

\ccsdesc[500]{Information systems~Information retrieval}
\ccsdesc[500]{Computing methodologies~Natural language processing}
\ccsdesc[300]{Human-centered computing~Visual analytics}

\keywords{GraphRAG, knowledge graphs, human-in-the-loop, skill compilation, provenance, curation agents, programming by demonstration}

\maketitle


\section{Introduction}
Graph-based Retrieval-Augmented Generation (GraphRAG) induces an entity--relation graph over a document collection and uses this structure to retrieve context for question answering~\cite{edge2025localglobalgraphrag}. Because the pipeline is not optimized end-to-end, local inaccuracies persist: entity fragmentation, spurious relations, and unstable clustering accumulate across stages and distort both graph structure and retrieved evidence, a phenomenon we call cascading error. In black-box settings these failures are hard to localize: answers stay fluent while the reasoning path is opaque.

Curation is the obvious remedy and the obvious bottleneck. No two induced graphs fail alike: domain, purpose, and language dictate which errors dominate, so every knowledge graph effectively demands its own cleaning toolchain. Highly inflective languages, for instance, fragment entities into inflected forms and abbreviations unless dedicated preprocessing absorbs them; elsewhere the dominant noise is spurious relations or plausible-but-wrong communities. Hard-coding a toolchain for one corpus does not transfer, and asking the curator to repeat the same edit hundreds of times does not scale.

We introduce GraphCraft, built on one stance: the journal of curation changes is training data. Two contributions follow. (i)~\textbf{Journal-to-skill compilation}: an agent induces the rule behind the curator's recent edits and a wizard materializes it at three price points, dry-run on the live graph and adopted only after preview; conflicting proposals queue for the user, and pinned nodes are off limits to every agent. (ii)~\textbf{End-to-end provenance}: every answer, suggestion, skill, and automatic invalidation traces back to the source spans and journal operations that produced it, so nothing in the graph or in an answer is unexplained. The remaining machinery, a bi-temporal store, a delta grammar rendered across time, edits, and queries, an activity timeline, incremental recomputation, multi-user editing, and pluggable reasoners, is presented as supporting tooling rather than as separate claims.

\section{Related Work}
XGraphRAG~\cite{wang2025xgraphraginteractivevisualanalysis} is the closest diagnostic system: visual analysis that traces suspicious retrievals through the pipeline, yet read-only over a fixed graph. AGENTiGraph~\cite{zhao2025agentigraph} translates chat intents into operations over manually curated KGs, whereas our agents learn from the curator's own journal and curate automatically induced artifacts. Zep/Graphiti~\cite{rasmussen2025zep} maintains a bi-temporal agent-memory graph with automatic fact invalidation but exposes no curation surface; our system makes its kind of structure, and its wrong invalidations, visible and human-correctable. OntoChat~\cite{zhang2024ontochat} explores conversational ontology engineering at requirement-elicitation time; our skills operate downstream, on the induced graph itself. iText2KG~\cite{lairgi2024itext2kg} and HippoRAG~\cite{gutierrez2024hipporag} build knowledge graphs incrementally but treat induction as a fully automated stage. General-purpose tools such as Neo4j Bloom~\cite{neo4j_bloom} and Cytoscape~\cite{shannon2003cytoscape} support persistent editing, but an edit there is a plain database write: nothing recomputes the summaries, communities, and embeddings that retrieval actually consumes, no provenance ties an edit to the answers it changes, and no mechanism turns repeated edits into automation.

\begin{table}[t]
\centering
\footnotesize
\setlength{\tabcolsep}{3.5pt}
\caption{Feature comparison across diagnostic and graph tools.}
\label{tab:feature_comparison}
\resizebox{\columnwidth}{!}{%
\begin{tabular}{lccccc}
\toprule
Feature & XGraphRAG & AGENTiGraph & Zep/Graphiti & Neo4j Bloom & Ours \\
\midrule
Pipeline diagnosis views        & \checkmark & \texttimes & \texttimes & \texttimes & \checkmark \\
Persistent edits affect retrieval & \texttimes & \checkmark & \texttimes & \checkmark & \checkmark \\
Ask $\rightarrow$ inspect $\rightarrow$ edit $\rightarrow$ re-ask & partial & partial & \texttimes & \texttimes & \checkmark \\
Curates induced (noisy) graphs & \checkmark & \texttimes & \texttimes & \texttimes & \checkmark \\
Bi-temporal model               & \texttimes & \texttimes & \checkmark & \texttimes & \checkmark \\
Change visualization (time/edit/query) & \texttimes & \texttimes & \texttimes & \texttimes & \checkmark \\
Human-correctable invalidation  & \texttimes & \texttimes & \texttimes & \texttimes & \checkmark \\
\bottomrule
\end{tabular}}%
\end{table}


\section{Skill Compilation}
\textbf{Journal-to-skill compilation.} The journal doubles as training data. The user selects their last $n$ operations; each is a structured witness that records the operation type, the nodes it touched with their types and layers, and the outcome. Rule induction runs three passes over these witnesses with the same pluggable LLM backend the pipeline uses for extraction: the first pass articulates the rule the curator appears to follow, the second applies that rule to the rest of the graph and proposes candidate operations, and the third verifies every candidate against the rule and grounds it in the live graph, dropping those whose targets no longer exist. The resulting skill is scoped to the whole graph, one layer, or one node type, and a wizard builds it at three price points: structural heuristics with no LLM calls, embedding candidates with one LLM verification pass, and the full three-pass reasoning above. Each build is dry-run on the live graph, its hypothetical effect rendered in the same delta grammar next to cost and latency badges, and the user adopts the cheapest build that does the job; accepted skills join the agent pool and carry their journal slice as provenance. Skills can contradict one another, one proposing to merge what another would retype. Proposals never auto-apply, and the user can pin a node after checking it by eye: a pinned node is excluded from every agent and skill, and competing proposals on unpinned nodes queue side by side for the user to settle.


\textbf{Why three price points.} The same rule can usually be enforced at very different costs. Morphological deduplication compiles to string and morphology heuristics that sweep the whole graph for free; relation cleanup may need embedding similarity with a single LLM verification pass; only genuinely contextual rules need the full three-pass reasoning. Exposing the ladder lets the user decide how much intelligence a recurring decision deserves, and the dry-run preview shows exactly what each tier would and would not catch before anything is applied.

\section{Provenance}
Provenance is enforced at every level rather than added per feature. Every node and edge carries spans back to the documents it was extracted from; selecting a chunk node shows its source text next to every entity extracted from it, so extraction is auditable in both directions. Every curation operation, whether issued by the user from the canvas or proposed by an agent, lands in a journal entry that records the actor, the affected set, and, for agent proposals, the rule and journal slice that generated them. Every answer is produced over an explicit evidence subgraph: the query is matched against node and chunk embeddings, expanded over the graph, and the LLM is prompted only over the attached chunk text, so each claim streams back with paragraph-level citations. And every automatic invalidation records which ingestion event caused it and why, so a wrong invalidation can be found and reverted, the failure mode a fully automatic backend cannot expose.

\section{Supporting Tools}
The compiler and the provenance contract sit on conventional but necessary machinery, which we list as tooling rather than claims. A bi-temporal store separates transaction time from event time, so an extraction error can be told apart from a genuine change in an evolving corpus. A delta grammar of born, dead, persisted, moved-community, and invalidated facts renders change identically across period diffs, edit cascades, and query evidence, and an activity timeline brushes any period. An incremental engine refreshes only the affected summaries, embeddings, and community assignments after an edit, with structural updates landing immediately and the rest completing in the background. Curation is multi-user in the manner of a collaborative editor: one variant is one shared document, operations apply in arrival order, and a teammate's edit ripples through the same delta grammar, signed with its author. Entity-type chips collapse an unreadable rendering into a legible map (Fig.~\ref{fig:project_setup}), a pluggable Think-on-Graph reasoner~\cite{sun2024thinkongraph} exposes multi-hop traversal as a heat overlay, and variants built under different configurations are compared in a synchronized split view. A self-contained batch mode runs without external services, and the full system is open source.

\begin{figure}[!tbp]
\centering
\begin{subfigure}{0.49\linewidth}
  \includegraphics[width=\linewidth]{graph_view_all.png}
  \caption{all node types}
\end{subfigure}\hfill
\begin{subfigure}{0.49\linewidth}
  \includegraphics[width=\linewidth]{graph_view_focus.png}
  \caption{people and organizations}
\end{subfigure}
\caption{The news graph with all node types (left) and with the type chips restricted to people and organizations (right).}
\label{fig:project_setup}
\end{figure}

\section{Demonstration}
We demonstrate the skill loop end to end; a recorded walkthrough is available at \url{https://vimeo.com/1141894235}.

\textbf{Case study: a morphology skill from twenty merges.} On a Russian podcast corpus, a highly inflective language, evidence highlighting reveals inflected variants and abbreviations of one department living as separate nodes across communities. The curator drills from a suspicious entity into its source chunk, merges duplicates by hand, and pins the nodes that are now verified. After a handful of such merges the curator selects them in the journal and compiles a skill: the induced rule reads, in effect, merge entities whose names are inflectional variants of the same lemma within one type. The cheap build, morphology rules with no LLM calls, reproduces nearly all manual merges on the dry run and proposes the remaining candidates across the corpus; the curator accepts it, and future ingestion batches are cleaned for free, with every proposal still arriving as a suggestion that traces back to the originating merges. On the evolving news corpus the same loop runs against weekly batches: the timeline shows what each week added, and accepted skills keep the new material consistent with the curated past.

\textbf{Human evaluation.} Six users worked with the system on the two corpora, issuing over 500 queries, applying edits, and re-running their questions; on a scored subset, answers were rated on factual correctness, coverage, and evidence alignment before and after curation (Table~\ref{tab:graphrag_results_human}). In a seeded error-correction task with six errors per graph and 15 minutes per participant, 78\% of seeded errors were corrected at 0.91 precision, fragmentation errors most reliably. Brief interactive curation restored graph fidelity without retraining.

\begin{table}[t]
\centering
\small
\begin{tabular*}{\columnwidth}{@{\extracolsep{\fill}}lccc}
\hline
\textbf{Dimension} & \textbf{Before} & \textbf{After} & $\Delta$ \\
\hline
Factual correctness & 0.67 & 0.72 & +0.05 \\
Coverage            & 0.63 & 0.69 & +0.06 \\
Evidence alignment  & 0.58 & 0.64 & +0.06 \\
\hline
\end{tabular*}
\caption{Human ratings before and after lightweight curation; 0--2 scores rescaled to $[0,1]$.}
\label{tab:graphrag_results_human}
\end{table}


\section{Limitations}
\textbf{Scalability.} The visualization comfortably handles graphs of up to 30{,}000 nodes, and with type chips and focused views a human rarely needs that many on screen at once. The practical bound today lies elsewhere: compiled skills start to take noticeably long once the active layer they run on exceeds roughly 800 nodes.
\textbf{Dependence on upstream extraction.} The system exposes and helps correct extraction noise but cannot recover information that was never extracted.
\textbf{Evaluation scope.} The evaluation is demo-oriented and validates qualitative behavior rather than large-scale statistical performance.

\section{Conclusion and Future Work}
We present GraphCraft, an interactive system that turns a curator's edit journal into reusable, cost-tiered skills and refuses to assert anything without provenance. Future work will arrange invariants, journal-induced rules, and LLM verification into an explicit hierarchy of oracles that spends human attention last, re-compile skills when the curator's behaviour drifts away from them, measure per build tier how many manual actions compiled skills save on held-out weeks, and add a cross-variant tier in which agreement between independently built graph variants auto-verifies a suggestion.

\section*{Ethical Statement}
All demonstration data consist of openly accessible podcast transcripts and news articles; no personal data are released with the system.

\bibliographystyle{ACM-Reference-Format}
\bibliography{sample-base}

\end{document}
